\documentclass[a4paper,11pt]{article}

\usepackage{graphicx}
\usepackage{booktabs}
\usepackage{amsmath}
\usepackage{hyperref}
\usepackage{caption}
\usepackage[margin=25mm]{geometry}

\graphicspath{{./}}
\hypersetup{colorlinks=true,linkcolor=blue,citecolor=blue,urlcolor=blue}

\title{\textbf{Effect of Thread Count and SIMD Vectorisation on Dense Matrix
Multiplication: an OpenMP Study on a Hybrid Multicore CPU}}
\author{<MY FULL NAME>\\Assignment 3}
\date{Big Data\\Grado en Ciencia e Ingeniería de Datos, ULPGC\\2026--2027}

\begin{document}
\maketitle

\begin{abstract}
This report studies the effect of OpenMP thread count and compiler
auto-vectorisation on dense double-precision matrix multiplication. Four
implementations were measured on one Windows 11 system: sequential scalar,
parallel scalar, sequential SIMD-enabled, and parallel SIMD-enabled. The
final clean session contains 1484 raw measurements and 66 measured
configurations. The $i$-$k$-$j$ kernel uses row-major matrices, and only the
multiplication is timed. At $n=1024$, parallel SIMD with $p=12$ has a
same-level speedup of 6.40 and efficiency of 0.533 relative to sequential
SIMD. At $n=4096$, the corresponding values are 2.93 and 0.244, based on
three repetitions. These values describe one hybrid CPU and one session; no
universal performance claim is made.
\end{abstract}

\section{Introduction and objective}

Dense matrix multiplication is a useful shared-memory experiment because its
computation is regular and its output rows can be distributed between
threads. This assignment investigates the following research question:
\emph{How does thread count affect the runtime, speedup, and efficiency of
dense matrix multiplication, and how much additional performance does SIMD
auto-vectorisation provide?}

The following statements are hypotheses, not results:
\begin{itemize}
  \item Speedup may increase approximately linearly up to the physical-core
  limit and then flatten or decrease when more logical threads are used.
  \item Small matrices may have lower efficiency because fixed overheads are
  significant.
  \item SIMD may improve both sequential and parallel implementations.
\end{itemize}

\section{Background}

OpenMP expresses parallel work over shared memory \cite{openmp}. In this experiment, the
outer matrix-row loop is distributed with a static schedule. Amdahl's law
describes the limit imposed by work that does not scale with the number of
threads \cite{amdahl}. The Roofline model relates attainable arithmetic
throughput to arithmetic intensity and memory bandwidth \cite{roofline}.

Let $T_{\mathrm{seq}}$ denote the median time of a matching sequential
baseline and let $T_p$ denote the median time with $p$ threads. Speedup and
efficiency are defined as
\[
 S(p)=\frac{T_{\mathrm{seq}}}{T_p},
 \qquad
 E(p)=\frac{S(p)}{p}.
\]
The statistic used for time is the median and variability is reported with
the interquartile range (IQR). Same-level speedup uses
\texttt{seq\_scalar} for \texttt{par\_scalar} and \texttt{seq\_simd} for
\texttt{par\_simd}. Combined speedup always uses \texttt{seq\_scalar}.

SIMD vectorisation evaluates independent operations in vector registers. The
$i$-$k$-$j$ loop order makes the innermost $j$ loop contiguous in memory.
This gives the compiler a regular loop to vectorise while keeping the
mathematical algorithm identical between builds.

\section{Methodology}

\subsection{Algorithm and variants}

The matrices are contiguous row-major arrays of \texttt{double}. The kernel
computes
\[
 C_{ij} \mathrel{+}= A_{ik}B_{kj}
\]
with loop order $i$-$k$-$j$. The input generator is xorshift64 with seed
\texttt{123456789}; generated values are uniform in $[0,1)$. Allocation,
generation, and clearing $C$ are outside the timed interval. A checksum is
computed after timing.

The exact builds are:
\begin{description}
\item[\texttt{seq\_scalar}] \texttt{-std=c11 -O3 -fno-tree-vectorize
-fno-tree-slp-vectorize -ffp-contract=off}.
\item[\texttt{par\_scalar}] The sequential scalar flags plus
\texttt{-fopenmp}.
\item[\texttt{seq\_simd}] \texttt{-std=c11 -O3 -march=native
-ffp-contract=off}.
\item[\texttt{par\_simd}] The sequential SIMD flags plus
\texttt{-fopenmp}.
\end{description}
No BLAS, NumPy, fast-math, or optimised library is used in the timed kernel.
The compiler is GCC 16.1.0, and the flag meanings are documented in the GCC
manual \cite{gcc}.

\subsection{Hardware and software}

The recorded system is Windows 11, an Intel Core i5-13420H, 8 physical cores,
12 logical processors, and 15.64 GB of RAM. Its topology is hybrid: 4
P-cores with Hyper-Threading and 4 E-cores. The compiler is MSYS2 UCRT64 GCC
16.1.0. Python is version 3.13.12; the recorded analysis environment also
contains NumPy 2.2.6, pandas 2.3.3, matplotlib 3.10.6, and psutil 7.1.0.
The power plan recorded in the system information was Equilibrado and AC
power was reported as true.

\subsection{Timing protocol}

On Windows, timing uses \texttt{QueryPerformanceCounter}; on other platforms
the source uses \texttt{clock\_gettime(CLOCK\_MONOTONIC)}. The timer includes
only the multiplication kernel. Input generation, allocation,
zero-initialisation, and checksum calculation are outside it.

The tested sizes are $n\in\{256,512,1024,2048,4096\}$. Sizes through
$n=1024$ use two warm-ups and 30 measured repetitions. The $n=2048$ cases
use one warm-up and 10 repetitions. The $n=4096$ cases use one warm-up and
three repetitions. For $n=256$ and $n=512$, ten multiplications are batched
and the CSV stores the average time per multiplication.

Parallel variants use $p\in\{1,2,4,6,8,12\}$ for the regular sizes, plus
$p=24$ at $n=1024$. Sequential variants use $p=1$. At $n=4096$, the
reduced coverage contains sequential scalar and SIMD, parallel SIMD at
$p\in\{1,4,8,12\}$, and parallel scalar at $p\in\{1,12\}$. A 120-second cap
applies to each multiplication.

The pinning test at $n=1024$ compared 30 repetitions without explicit CPU
affinity against 30 repetitions pinned to logical CPU 0. Unpinned
\texttt{seq\_simd} had median 0.285 s and IQR 0.0362 s; CPU-0-pinned
\texttt{seq\_simd} had median 0.260 s and IQR 0.0169 s. The final session
therefore pins sequential and $p=1$ runs to CPU 0; runs with $p>1$ are not
pinned. The environment sets \texttt{OMP\_NUM\_THREADS},
\texttt{OMP\_PROC\_BIND=close}, and \texttt{OMP\_PLACES=cores}. MinGW
\texttt{libgomp} reported that affinity was not supported and displayed an
empty \texttt{OMP\_PLACES}.

\begin{table}[htbp]
\centering
\caption{Configuration sets represented in the final clean session.}
\begin{tabular}{@{}llll@{}}
\toprule
Size & Sequential variants & Parallel thread counts & Repetitions\\
\midrule
256, 512, 1024 & scalar and SIMD, $p=1$ &
$p=1,2,4,6,8,12$; $p=24$ at 1024 & 30\\
2048 & scalar and SIMD, $p=1$ &
$p=1,2,4,6,8,12$ & 10\\
4096 & scalar and SIMD, $p=1$ &
SIMD $p=1,4,8,12$; scalar $p=1,12$ & 3\\
\bottomrule
\end{tabular}
\end{table}

\subsection{Correctness and memory}

Correctness tests cover identity, zero, a hand-computed $3\times3$ product,
and NumPy reference results for $n=64$, $n=128$, and $n=256$. The comparison
uses $\mathrm{atol}=\mathrm{rtol}=10^{-9}$. Sequential and parallel result
checksums match exactly in the tested configurations. Each thread writes
only its assigned rows of $C$, so the outer parallel loop has no write race.

On Windows, peak memory uses \texttt{GetProcessMemoryInfo} and
\texttt{PeakWorkingSetSize}. The theoretical array-storage metric is
$3\cdot8\cdot n^2$ bytes. These are different quantities. The memory budget
was 2 GB.

\setcounter{subsection}{5}
\subsection{Vectorisation evidence}

The saved GCC report for \texttt{seq\_simd} states: ``\texttt{matmul.c:57:26:
optimized: loop vectorized using 16 byte vectors and unroll factor 2}''.
The saved \texttt{par\_simd} disassembly contains the following hot-loop
excerpt:
\begin{verbatim}
vmulpd (%rdx,%r8,1),%ymm1,%ymm0
vaddpd (%rax,%r8,1),%ymm0,%ymm0
vmovupd %ymm0,(%rax,%r8,1)
\end{verbatim}
The vectoriser report also records 32-byte vectors in the SIMD build. The
assembly uses AVX instructions with YMM registers; the evidence files do not
identify a narrower or wider named SIMD extension beyond these instructions.

\subsection{Discarded runs}

The first discarded run had the laptop lid closed while a game was running.
The second used the coarse OpenMP timer observed in the parallel builds. The
third preliminary run applied the 120-second limit to the whole process
instead of to each multiplication. All are retained under
\texttt{results/archive/} and excluded from the final analysis.

\section{Results}

All values in this section come from the final \texttt{summary.csv},
\texttt{raw\_measurements.csv}, or \texttt{sanity\_report.md}. Times are
seconds per multiplication. Error values are IQR and $r$ is the number of
measured repetitions. Sentences in this section are descriptive only.

\subsection{Runtime}

\begin{table}[htbp]
\centering
\scriptsize
\caption{Median time in seconds for every tested thread count. A dash means
that the configuration was not part of the final session. Values are rounded
to three significant digits.}
\begin{tabular}{@{}llrrrrrrr@{}}
\toprule
Variant & $n$ & $p=1$ & $p=2$ & $p=4$ & $p=6$ & $p=8$ & $p=12$ & $p=24$\\
\midrule
\texttt{seq\_scalar} & 256 & 0.00737 & -- & -- & -- & -- & -- & --\\
\texttt{seq\_scalar} & 512 & 0.0585 & -- & -- & -- & -- & -- & --\\
\texttt{seq\_scalar} & 1024 & 0.501 & -- & -- & -- & -- & -- & --\\
\texttt{seq\_scalar} & 2048 & 5.36 & -- & -- & -- & -- & -- & --\\
\texttt{seq\_scalar} & 4096 & 48.7 & -- & -- & -- & -- & -- & --\\
\texttt{seq\_simd} & 256 & 0.00226 & -- & -- & -- & -- & -- & --\\
\texttt{seq\_simd} & 512 & 0.0275 & -- & -- & -- & -- & -- & --\\
\texttt{seq\_simd} & 1024 & 0.262 & -- & -- & -- & -- & -- & --\\
\texttt{seq\_simd} & 2048 & 3.77 & -- & -- & -- & -- & -- & --\\
\texttt{seq\_simd} & 4096 & 33.2 & -- & -- & -- & -- & -- & --\\
\texttt{par\_scalar} & 256 & 0.00762 & 0.00420 & 0.00194 & 0.00184 & 0.00150 & 0.00123 & --\\
\texttt{par\_scalar} & 512 & 0.0581 & 0.0296 & 0.0144 & 0.0140 & 0.0113 & 0.00916 & --\\
\texttt{par\_scalar} & 1024 & 0.503 & 0.268 & 0.145 & 0.123 & 0.102 & 0.0749 & 0.0776\\
\texttt{par\_scalar} & 2048 & 5.34 & 3.30 & 1.77 & 1.46 & 1.18 & 0.964 & --\\
\texttt{par\_scalar} & 4096 & 47.0 & -- & -- & -- & -- & 12.3 & --\\
\texttt{par\_simd} & 256 & 0.00223 & 0.00117 & 0.000589 & 0.000638 & 0.000724 & 0.000646 & --\\
\texttt{par\_simd} & 512 & 0.0275 & 0.0122 & 0.00665 & 0.00561 & 0.00514 & 0.00397 & --\\
\texttt{par\_simd} & 1024 & 0.268 & 0.127 & 0.0635 & 0.0614 & 0.0486 & 0.0409 & 0.0404\\
\texttt{par\_simd} & 2048 & 3.78 & 2.20 & 1.17 & 0.970 & 0.901 & 0.767 & --\\
\texttt{par\_simd} & 4096 & 33.3 & -- & 10.9 & -- & 11.5 & 11.3 & --\\
\bottomrule
\end{tabular}
\end{table}

\begin{figure}[htbp]
\centering
\includegraphics[width=\linewidth]{time_vs_n.png}
\caption{Median time per multiplication versus $n$ on log--log axes for
$p=1$, $p=4$, and $p=12$; error bars are IQR.}
\label{fig:time}
\end{figure}

\subsection{Speedup and efficiency}

\begin{table}[htbp]
\centering
\scriptsize
\caption{Same-level speedup/efficiency for every tested thread count,
shown as $S(p)/E(p)$. Scalar rows use \texttt{seq\_scalar}; SIMD rows use
\texttt{seq\_simd}. Values are rounded to three significant digits.}
\begin{tabular}{@{}llrrrrrrr@{}}
\toprule
Variant & $n$ & $p=1$ & $p=2$ & $p=4$ & $p=6$ & $p=8$ & $p=12$ & $p=24$\\
\midrule
\texttt{par\_scalar} & 256 & 0.966/0.966 & 1.75/0.876 & 3.80/0.951 & 4.01/0.668 & 4.92/0.615 & 6.00/0.500 & --\\
\texttt{par\_scalar} & 512 & 1.01/1.01 & 1.98/0.989 & 4.05/1.01 & 4.18/0.696 & 5.20/0.650 & 6.39/0.532 & --\\
\texttt{par\_scalar} & 1024 & 0.995/0.995 & 1.87/0.934 & 3.45/0.863 & 4.07/0.678 & 4.93/0.616 & 6.69/0.557 & 6.45/0.269\\
\texttt{par\_scalar} & 2048 & 1.00/1.00 & 1.63/0.813 & 3.03/0.758 & 3.66/0.611 & 4.52/0.565 & 5.56/0.463 & --\\
\texttt{par\_scalar} & 4096 & 1.04/1.04 & -- & -- & -- & -- & 3.96/0.330 & --\\
\texttt{par\_simd} & 256 & 1.01/1.01 & 1.93/0.967 & 3.84/0.959 & 3.54/0.590 & 3.12/0.390 & 3.50/0.291 & --\\
\texttt{par\_simd} & 512 & 1.00/1.00 & 2.26/1.13 & 4.14/1.04 & 4.91/0.818 & 5.36/0.670 & 6.94/0.578 & --\\
\texttt{par\_simd} & 1024 & 0.977/0.977 & 2.06/1.03 & 4.12/1.03 & 4.26/0.710 & 5.38/0.673 & 6.40/0.533 & 6.47/0.270\\
\texttt{par\_simd} & 2048 & 0.997/0.997 & 1.71/0.855 & 3.22/0.806 & 3.88/0.647 & 4.18/0.522 & 4.91/0.409 & --\\
\texttt{par\_simd} & 4096 & 0.998/0.998 & -- & 3.06/0.765 & -- & 2.88/0.360 & 2.93/0.244 & --\\
\bottomrule
\end{tabular}
\end{table}

\begin{figure}[htbp]
\centering
\includegraphics[width=\linewidth]{speedup_vs_threads.png}
\caption{Same-level speedup versus $p$ in two rows of size panels. The
baseline is sequential scalar for scalar and sequential SIMD for SIMD. The
ideal line and the physical and logical processor markers are shown; error
bars represent IQR-derived uncertainty.}
\end{figure}

\begin{figure}[htbp]
\centering
\includegraphics[width=\linewidth]{efficiency_vs_threads.png}
\caption{Same-level efficiency versus $p$ in two rows of size panels, using
the same baselines as the speedup figure. Error bars represent IQR-derived
uncertainty.}
\end{figure}

\subsection{OpenMP $p=1$ overhead}

\begin{table}[htbp]
\centering
\caption{OpenMP $p=1$ speedup and efficiency relative to the matching
sequential baseline.}
\begin{tabular}{@{}lrrrrr@{}}
\toprule
Variant & $n$ & $S(1)$ & $E(1)$ & IQR (s) & $r$\\
\midrule
\texttt{par\_scalar} & 256 & 0.966 & 0.966 & 0.000106 & 30\\
\texttt{par\_simd} & 256 & 1.01 & 1.01 & 0.000140 & 30\\
\texttt{par\_scalar} & 512 & 1.01 & 1.01 & 0.00160 & 30\\
\texttt{par\_simd} & 512 & 1.00 & 1.00 & 0.000970 & 30\\
\texttt{par\_scalar} & 1024 & 0.995 & 0.995 & 0.0235 & 30\\
\texttt{par\_simd} & 1024 & 0.977 & 0.977 & 0.0186 & 30\\
\texttt{par\_scalar} & 2048 & 1.00 & 1.00 & 0.110 & 10\\
\texttt{par\_simd} & 2048 & 0.997 & 0.997 & 0.0434 & 10\\
\texttt{par\_scalar} & 4096 & 1.04 & 1.04 & 1.52 & 3\\
\texttt{par\_simd} & 4096 & 0.998 & 0.998 & 1.21 & 3\\
\bottomrule
\end{tabular}
\end{table}

\begin{figure}[htbp]
\centering
\includegraphics[width=.85\linewidth]{openmp_overhead.png}
\caption{Runtime of parallel variants at $p=1$ and their matching
sequential baselines. Error bars are IQR.}
\end{figure}

\subsection{Vectorisation and GFLOP/s}

\begin{table}[htbp]
\centering
\caption{Two-by-two vectorisation view at $n=2048$ and $n=4096$.}
\begin{tabular}{@{}lrrrrr@{}}
\toprule
Variant & $n$ & $p$ & Median (s) & IQR (s) & $r$\\
\midrule
\texttt{seq\_scalar} & 2048 & 1 & 5.36 & 0.185 & 10\\
\texttt{seq\_simd} & 2048 & 1 & 3.77 & 0.0619 & 10\\
\texttt{par\_scalar} & 2048 & 12 & 0.964 & 0.0402 & 10\\
\texttt{par\_simd} & 2048 & 12 & 0.767 & 0.136 & 10\\
\texttt{seq\_scalar} & 4096 & 1 & 48.7 & 0.957 & 3\\
\texttt{seq\_simd} & 4096 & 1 & 33.2 & 1.29 & 3\\
\texttt{par\_scalar} & 4096 & 12 & 12.3 & 0.683 & 3\\
\texttt{par\_simd} & 4096 & 12 & 11.3 & 1.31 & 3\\
\bottomrule
\end{tabular}
\end{table}

\begin{figure}[htbp]
\centering
\includegraphics[width=.85\linewidth]{vectorisation_2x2.png}
\caption{Scalar/SIMD and one/12-thread runtimes at $n=2048$ and $n=4096$.
Bars are medians and error bars are IQR.}
\end{figure}

\begin{table}[htbp]
\centering
\caption{Sequential GFLOP/s calculated as $2n^3/T_{\mathrm{seq}}/10^9$.}
\begin{tabular}{@{}lrrrr@{}}
\toprule
Variant & $n$ & GFLOP/s & IQR (s) & $r$\\
\midrule
\texttt{seq\_scalar} & 256 & 4.56 & 0.000651 & 30\\
\texttt{seq\_simd} & 256 & 14.9 & 0.000314 & 30\\
\texttt{seq\_scalar} & 512 & 4.59 & 0.00249 & 30\\
\texttt{seq\_simd} & 512 & 9.75 & 0.000780 & 30\\
\texttt{seq\_scalar} & 1024 & 4.29 & 0.0260 & 30\\
\texttt{seq\_simd} & 1024 & 8.21 & 0.0187 & 30\\
\texttt{seq\_scalar} & 2048 & 3.21 & 0.185 & 10\\
\texttt{seq\_simd} & 2048 & 4.56 & 0.0619 & 10\\
\texttt{seq\_scalar} & 4096 & 2.82 & 0.957 & 3\\
\texttt{seq\_simd} & 4096 & 4.14 & 1.29 & 3\\
\bottomrule
\end{tabular}
\end{table}

\begin{figure}[htbp]
\centering
\includegraphics[width=.8\linewidth]{gflops_vs_n.png}
\caption{Sequential GFLOP/s versus $n$, calculated from median time. Error
bars represent propagated IQR.}
\end{figure}

\subsection{Memory, oversubscription, and $n=4096$}

\begin{table}[htbp]
\centering
\caption{Peak working set measured with Windows
\texttt{PeakWorkingSetSize}, compared with theoretical array storage. For
$n=256$ and $n=512$, batch mode allocates one separate result buffer per
batched multiplication, so the measured process peak includes these buffers
and is not the storage of only three arrays.}
\begin{tabular}{@{}rrrr@{}}
\toprule
$n$ & Peak median (MiB) & Peak IQR (MiB) & Theoretical (MiB)\\
\midrule
256 & 9.71 & 0 & 1.50\\
512 & 27.7 & 0 & 6.00\\
1024 & 27.7 & 0 & 24.0\\
2048 & 99.7 & 0 & 96.0\\
4096 & 388 & 0.102 & 384\\
\bottomrule
\end{tabular}
\end{table}

\begin{figure}[htbp]
\centering
\includegraphics[width=.85\linewidth]{peak_memory_vs_n.png}
\caption{Measured peak working set and theoretical three-array storage.
Measured points are medians with IQR error bars.}
\end{figure}

\begin{table}[htbp]
\centering
\caption{Oversubscription comparison at $n=1024$. The ratio is
$T_{12}/T_{24}$; $p=24$ is the oversubscribed configuration.}
\begin{tabular}{@{}lrrrr@{}}
\toprule
Variant & $T_{12}/T_{24}$ & $T_{12}$ (s) & $T_{24}$ (s) & $r$\\
\midrule
\texttt{par\_scalar} & 0.965 & 0.0749 & 0.0776 & 30\\
\texttt{par\_simd} & 1.01 & 0.0409 & 0.0404 & 30\\
\bottomrule
\end{tabular}
\end{table}

\begin{figure}[htbp]
\centering
\includegraphics[width=.75\linewidth]{oversubscription_1024.png}
\caption{Relative $p=12$ versus $p=24$ speed at $n=1024$ for the parallel
variants. Error bars are based on IQR.}
\end{figure}

The $n=4096$ measurements use three repetitions and the reduced configuration
set stated in the methodology. The sequential scalar median is 48.7 s, the
sequential SIMD median is 33.2 s, and the parallel SIMD $p=12$ median is
11.3 s.

\section{Discussion}

\textbf{TODO (student).} No interpretation is supplied here.
\begin{itemize}
  \item Where does speedup stop increasing, and how might the hybrid topology
  relate to that point?
  \item What does the $p=1$ overhead show for scalar and SIMD variants?
  \item What does the 2x2 vectorisation view show at the two largest sizes?
  \item Do the GFLOP/s and memory figures support a cache or memory-bound
  explanation?
  \item How should the IQR and sanity-report flags affect confidence?
\end{itemize}

\subsection{Limitations}

The processor is hybrid, with 4 P-cores with Hyper-Threading and 4 E-cores;
thread counts above the tested homogeneous subset mix core types. Sequential
and $p=1$ runs are pinned to CPU 0, while runs with $p>1$ are not pinned.
MinGW \texttt{libgomp} reported that \texttt{OMP\_PROC\_BIND} and
\texttt{OMP\_PLACES} could not enforce affinity. The $n=4096$ set is reduced
and has three repetitions. The experiment uses one machine and one session;
thermal state and background activity were not fully controlled. The kernel
is not cache-blocked, so memory-bound behaviour was not tested directly.
Peak working set and theoretical array storage are different metrics.

\section{Conclusions}

\textbf{TODO (student).}
\begin{itemize}
  \item Which hypotheses are supported or rejected by the measured data?
  \item What is the measured contribution of SIMD relative to threading?
  \item Which thread counts are most defensible for this hybrid CPU?
  \item Which conclusions remain limited by the protocol and reduced
  $n=4096$ coverage?
\end{itemize}

\section{AI usage disclosure}

An AI coding agent (Google Antigravity) was used for code generation,
debugging, benchmark-driver development, analysis scripts, figure generation,
and report scaffolding. I validated the correctness tests and measurements
and am responsible for the final content, interpretation, and submission.

\section{References}

\begin{thebibliography}{9}

\bibitem{amdahl}
G.~M. Amdahl, ``Validity of the single processor approach to achieving large
scale computing capabilities,'' in \emph{AFIPS Spring Joint Computer
Conference}, 1967, pp. 483--485.

\bibitem{openmp}
OpenMP Architecture Review Board, \emph{OpenMP Application Programming
Interface, Version 5.2}, Nov. 2021.

\bibitem{roofline}
S.~Williams, A.~Waterman, and D.~Patterson, ``Roofline: an insightful visual
performance model for multicore architectures,'' \emph{Communications of the
ACM}, vol. 52, no. 4, pp. 65--76, 2009.

\bibitem{gcc}
Free Software Foundation, \emph{Using the GNU Compiler Collection (GCC), for
GCC version 16.1.0}, GCC documentation, 2026. Available:
\url{https://gcc.gnu.org/onlinedocs/gcc-16.1.0/gcc/}

\end{thebibliography}

\noindent\textbf{Repository:} \texttt{<MY\_GITHUB\_URL>}. The reproducibility
materials are the C source, Python driver and analysis scripts, final raw CSV,
summary CSV, figures, vectorisation evidence, README, and this report.

\end{document}
